\subsection{The relative $S_3$ diagram for \S47}

The additional condition of \S47 concerns two different homomorphisms.
Its anchors are cote 140--3, folios 415--420, scans 229--240, especially
equations (14), (15), (20), and (30). We use the explicitly checked
dictionary in the supplied LM2 return, B.m3 and E.4, and its warning
B.g12 about finite quotients. The original cote-140--3 images were not
available for a fresh image check in this return. The definitions and
proofs below are independent algebraic statements about the displayed
maps; they do not depend on a doubtful manuscript letter.

Put
\[
 B=\langle s,r\mid s^2=r^3=1\rangle,
 \qquad p:B\longrightarrow S_3,
 \qquad K=\ker p,
\]
where $p(s)$ is a transposition and $p(r)$ a three-cycle. The group $K$
is free on $x=(rs)^2$ and $y=(sr)^2$; set $z=(xy)^{-1}$. Distinguish
the abstract free group $F_0=\langle X,Y\rangle$, $Z=(XY)^{-1}$, and
the two maps
\[
 \begin{aligned}
 \iota:F_0&\stackrel{\sim}{\longrightarrow}K\subset B,
      &(X,Y)&\longmapsto(x,y),\\
 \varphi:F_0&\twoheadrightarrow B,
      &(X,Y,Z)&\longmapsto(s,sr,r^{-1}).
 \end{aligned}
\]
Thus $v=p\varphi$ is a second quotient onto $S_3$, and
$N=\ker v$ is a different free group. Write $c:N\hookrightarrow F_0$
for inclusion and $d=\varphi|_N:N\twoheadrightarrow K$ for restriction.
All Lie algebras in this subsection are over $\Q$; degree $n$ means
that a derivation raises bracket degree by $n$.

\begin{statement}{B.25}{Proved}{The correct relative source has rank seven}
Relative completion with respect to the indicated finite quotients has
exact sequences
\[
 \begin{aligned}
 1\longrightarrow U(K)\longrightarrow\mathscr R_p(B)
       &\longrightarrow S_3\longrightarrow1,\\
 1\longrightarrow U(N)\longrightarrow\mathscr R_v(F_0)
       &\longrightarrow S_3\longrightarrow1,
 \end{aligned}
\]
where $U(K)$ and $U(N)$ are free prounipotent groups of ranks two and
seven. The quotient $\varphi$ induces
$\Phi:\mathscr R_v(F_0)\to\mathscr R_p(B)$, the identity on $S_3$,
whose map on radicals is $U(d)$. The absolute completion map induces
$\Theta:\mathscr R_v(F_0)\to U(F_0)$, whose map on radicals is $U(c)$.
The isomorphism $\iota_U:U(F_0)\to U(K)$ is a third map.

The associated graded source is the Lie algebra of an eight-punctured
sphere:
\[
 L_7=L(A_0,A_1,A_2,B_0,B_1,B_2,C_0,C_1)
       /(A_0+A_1+A_2+B_0+B_1+B_2+C_0+C_1).
\]
Writing both rank-two targets as $L_2=L(X,Y,Z)/(X+Y+Z)$ through
$\iota_U$, the maps are
\[
 \begin{array}{c|ccc}
       & A_i&B_i&C_i\\ \hline
 c     &2X&2Y&3Z\\
 d     &0&(X,Y,Z)_i&0.
 \end{array}
\]
The deck group permutes each three-element cluster $A$ and $B$ by
the standard $S_3$ action and permutes the two $C$ punctures by sign.
It acts trivially on the base target of $c$ and by puncture permutation
on the target of $d$. In degree one the joint map $(c,d)$ has rank four
and kernel isomorphic to $\mathrm{sgn}\oplus\mathrm{std}$.
\end{statement}
\begin{proof}
Relative completion means the universal prounipotent extension of the
chosen reductive quotient admitting the given homomorphism of the
discrete group; a finite group in characteristic zero is allowed
\citep[\S1, pp.~47--48]{Hain1998Relative}. For a finite quotient with
kernel $H$, its radical is $U(H)$. To see the relevant exactness,
restrict any representation of the relative completion to $H$.
Conversely, induce a unipotent representation of $H$ over the finitely
many cosets. Its algebraic image has the prescribed finite permutation
quotient and a subgroup of a product of unipotent groups as kernel.
This shows that every unipotent quotient of $H$ is detected by the
relative completion. Restriction and induction prove the universal
identification, as in B.8. Naturality gives $\Phi$ and $\Theta$.
Schreier's formula gives $\operatorname{rank}N=1+6(2-1)=7$.
The orbifold Euler characteristic of $C_2*C_3$ is $-1/6$, so its
torsion-free index-six kernel $K$ has Euler characteristic $-1$ and
rank two; Reidemeister rewriting gives the displayed basis $x,y$.

The maps can also be seen simultaneously in the rational covering
\[
 T(\lambda)=
 \frac{27\lambda^2(1-\lambda)^2}{4(1-\lambda+\lambda^2)^3}.
\]
Remove $0,1,\infty,-1,1/2,2,\zeta,\bar\zeta$ from the source,
where $\zeta=(1+\sqrt{-3})/2$, and remove $0,1,\infty$ from the
$T$-line. This is the regular degree-six cover with monodromy $v$.
The deck transformations $\lambda\mapsto1-\lambda$ and
$\lambda\mapsto1/(1-\lambda)$ generate $S_3$.
The punctures over $T=1,0,\infty$ have local degrees
$(2,2,2)$, $(2,2,2)$, and $(3,3)$, respectively. In particular,
\[
 4(1-\lambda+\lambda^2)^3-27\lambda^2(1-\lambda)^2
  =(\lambda+1)^2(2\lambda-1)^2(\lambda-2)^2.
\]
Euler characteristic gives genus zero and eight punctures.
Label the punctures over $1$ by
$(A_0,A_1,A_2)=(-1,2,1/2)$, and those over $0$ by
$(B_0,B_1,B_2)=(1,0,\infty)$. The two roots of
$1-\lambda+\lambda^2$ form the $C$ cluster. Filling the five
$A$ and $C$ punctures gives the target tripod in the $\lambda$-line.
It realizes $d$; the covering realizes $c$. Peripheral classes in
homology therefore have the displayed images. Since all three graded
Lie algebras are free on their first graded pieces, these first-degree
formulas specify the complete graded Lie homomorphisms.

The permutation representation on eight punctures has character
$(8,2,2)$ on the identity, a transposition, and a three-cycle.
Removing its diagonal line gives
\[
 H_1(N,\Q)\cong2\mathbf1\oplus\mathrm{sgn}\oplus2\mathrm{std},
 \qquad\chi=(7,1,1).
\]
In the free basis $A_0,A_1,A_2,B_0,B_1,B_2,C_0$, the two matrices are
\[
 c=\begin{pmatrix}2&2&2&0&0&0&-3\\0&0&0&2&2&2&-3\end{pmatrix},
 \qquad
 d=\begin{pmatrix}0&0&0&1&0&-1&0\\0&0&0&0&1&-1&0\end{pmatrix}.
\]
They have joint rank four, proving the last assertion.
\end{proof}

Here is a correctly typed formulation of the extra datum. For specified
automorphisms $a_v$ of $\mathscr R_v(F_0)$, $a$ of $U(F_0)$, and $b$
of $\mathscr R_p(B)$, with $a_v,b$ inducing the identity on the finite
quotients, the strict diagram asks for
\[
 \Theta a_v=a\Theta,\qquad \Phi a_v=b\Phi,\qquad
 b|_{U(K)}=\iota_Ua\iota_U^{-1}.
\]
An outer diagram allows the appropriate target inner automorphisms.
These are not equations involving $a$ alone. The absolute automorphism
$a$ does not specify $a_v$ on a rank-seven radical. In particular,
writing $U_2\to\mathscr R_p(B)$ in place of $\Phi$ would discard its
nontrivial component map and change the problem.

\begin{statement}{B.26}{Proved}{Every invariant graded action has a compatible relative lift}
Let $\mathfrak s_{r,n}$ denote degree-$n$ derivations tangential to
every peripheral class of the $(r+1)$-punctured sphere, modulo inner
derivations, as in B.10. For each $n\ge1$ the maps in B.25 induce a
surjection
\[
 R_n:\mathfrak s_{7,n}^{S_3}
       \longrightarrow
       \mathfrak s_{2,n}\oplus\mathfrak s_{2,n}^{S_3},
 \qquad [D]\longmapsto(c_*[D],d_*[D]).
\]
Consequently, for every $q\in\mathfrak s_{2,n}^{S_3}$ there is a
peripheral, deck-invariant relative graded lift whose two induced
outer actions are $(q,q)$. If $q\ne0$, there is also a lift with the
same absolute projection and induced pair $(q,0)$, which fails the
compatibility condition. Thus existential graded compatibility
imposes no additional equation on the special invariant $q$.
\end{statement}
\begin{proof}
First the map is defined. An invariant outer source class has an
invariant representative by averaging over the finite deck group.
A tangential derivation preserves the ideal generated by the $A$ and
$C$ punctures, which is the kernel of $d$. The map $c$ is the maximal
graded quotient on which the deck group acts trivially: its kernel
is the ideal generated by $gv-v$ in degree one. Indeed the
coinvariants of $H_1(N,\Q)$ are two-dimensional and $c$ is their
isomorphism onto $\Q X\oplus\Q Y$. An equivariant derivation
preserves that ideal. Its two descendants are tangential, and the
second is invariant. Changing the source representative by an inner
derivation changes both descendants by inner derivations, so $R_n$
is well-defined on outer classes.

We construct independent lifts for both factors. Write a representative
of $a\in\mathfrak s_{2,n}$ as
\[
 a(e)=[e,a_e(X,Y)],\quad e=X,Y,Z,\qquad
 \sum_{e=X,Y,Z}[e,a_e]=0.
\]
Set
\[
 X'=\tfrac16\sum_iA_i,\qquad
 Y'=\tfrac16\sum_iB_i,\qquad
 Z'=\tfrac16\sum_iC_i=-X'-Y'.
\]
Define its cabling $C(a)$ by
\[
 \begin{aligned}
 C(a)(A_i)&=[A_i,a_X(X',Y')],\\
 C(a)(B_i)&=[B_i,a_Y(X',Y')],\\
 C(a)(C_i)&=[C_i,a_Z(X',Y')].
 \end{aligned}
\]
The image of the peripheral relation is
$6\sum_e[e',a_e(X',Y')]=0$. Thus this is a tangential derivation of
$L_7$. It is deck-invariant, since every cluster sum is fixed.
Moreover $c(X')=X$, $c(Y')=Y$, and $d(X')=d(Y')=0$. Positive-degree
conjugators vanish at zero, so
\[
                 R_n(C(a))=(a,0).
\]
A common inner change of the three $a_e$ produces a source inner
change, so the construction also respects outer classes.

For the second factor take $b\in\mathfrak s_{2,n}^{S_3}$. The
rank-two outer space is zero in degree one, so assume $n\ge2$.
Average a representative of $b$ over $S_3$. In these degrees the
centralizer of a peripheral generator in the free Lie algebra has
no degree-$n$ part; hence its three conjugators $b_X,b_Y,b_Z$ are
unique and equivariant. Subtract their common mean, an inner change,
to arrange
\[
 b_X+b_Y+b_Z=0,\qquad g(b_e)=b_{ge}.
\]
Subtracting the common conjugator $b_Z$ gives a representative $b^0$
with conjugators $(b_X-b_Z,b_Y-b_Z,0)$. It fixes $Z$ and therefore
fixes $X+Y$. Substitute $X=B_0$, $Y=B_1$ in its first two images,
and set the images of $A_0,A_1,A_2,B_2,C_0,C_1$ to zero. Call the
result $E$. The identity $b^0(X+Y)=0$ proves that $E$ preserves the
eight-puncture relation; it is tangential to all eight classes.
Its $d$-descendant is $b^0$. Its $c$-descendant is zero, because
$c(B_0)=c(B_1)=2Y$ and a homogeneous Lie polynomial of degree at
least two in one letter is zero.

Finally put $T(b)=\frac16\sum_g gEg^{-1}$. It is invariant and
still has zero $c$-descendant. Its $d$-descendant is the average of
$b^0$. The averaged inner discrepancy is the inner derivation with
conjugator
$\frac16\sum_g g(b_Z)=\frac13(b_X+b_Y+b_Z)=0$.
Consequently
\[
                 R_n(T(b))=(0,b).
\]
Combining the two constructions proves surjectivity in all positive
degrees. The classes $C(q)+T(q)$ and $C(q)$ have the two pairs asserted
in the statement. When $q$ is nonzero as an outer class, their
different target projections remain different modulo inner derivations.
\end{proof}

\begin{statement}{B.27}{Proved}{Exact relative dimensions through weight nine}
Put $s_n=\dim\mathfrak s_{7,n}^{S_3}$,
$a_n=\dim\mathfrak s_{2,n}$, and
$t_n=\dim\mathfrak s_{2,n}^{S_3}$. Write $k_n=\dim\ker R_n$ and
$h_n=\dim R_n^{-1}(\operatorname{diag}\mathfrak s_{2,n}^{S_3})$.
The following dimensions are exact:
\[
\begin{array}{c|rrrrrrrrr}
 n&1&2&3&4&5&6&7&8&9\\ \hline
 s_n&4&5&38&120&693&3163&17356&92930&524166\\
 a_n&0&0&1&0&3&0&6&4&13\\
 t_n&0&0&1&0&1&0&2&1&4\\
 k_n&4&5&36&120&689&3163&17348&92925&524149\\
 h_n&4&5&37&120&690&3163&17350&92926&524153
\end{array}
\]
The image of the compatible relative space in the absolute special
invariants has dimension $t_n$. In particular the explicit $q_7$ of
B.12 passes this existential relative graded test.
\end{statement}
\begin{proof}
Let $V=H_1(N,\Q)$ and let $P$ be the permutation module on its eight
punctures. Their characters on $(1,(12),(123))$ are $(7,1,1)$ and
$(8,2,2)$. The exact special-derivation sequence in B.10 gives
\[
 \operatorname{tr}(g\mid\mathfrak s_{7,n})
 =\chi_V(g)\operatorname{tr}(g\mid L_nV)
   -\operatorname{tr}(g\mid L_{n+1}V)-\delta_{n,1}\chi_P(g),
\]
where
\[
 \operatorname{tr}(g\mid L_nV)
   =\frac1n\sum_{d\mid n}\mu(d)\chi_V(g^d)^{n/d}.
\]
Averaging over the three conjugacy classes, with sizes $1,3,2$,
gives $s_n$. The analogous rank-two formula gives $a_n$ and $t_n$.
For example the source characters in weights seven, eight, and nine
are $(103236,300,0)$, $(558404,-300,38)$, and
$(3140032,1680,-38)$. These give $17356$, $92930$, and $524166$.
Surjectivity in B.26 gives
\[
 \dim\ker R_n=s_n-a_n-t_n,\qquad
 \dim R_n^{-1}(\operatorname{diag}\mathfrak s_{2,n}^{S_3})=s_n-a_n.
\]

The executable \texttt{b\_relative47.py} independently recomputes
every rank-two special invariant basis through weight nine by rational
row reduction in standard Lyndon brackets. It verifies antisymmetry,
$J=0$, the averaged zero-sum equivariant conjugators, and both
projections of the two lifts for every basis vector. It also applies
the construction to the specified integral vector $q_7$, with its
actual coefficients from B.12. Rational associative-word certificates
for all these conjugators are saved in \texttt{b\_relative47.json};
the dimension table and executed output are saved as CSV and
\texttt{.out}. This exhausts the indicated rational rank-two spaces,
not merely selected vectors. The much larger rank-seven dimensions
come from the exact character identities just proved; no expanded
rank-seven matrix or sampled modular rank is claimed. The construction
in B.26 supplies the source tangentiality identities in all weights.
\end{proof}

\begin{statement}{B.28}{Open}{The prescribed profinite word remains a separate test}
The preceding results do not determine whether a prescribed profinite
loop automorphism satisfies the normalized \S47 equation
\[
 \xi=\bigl(\varphi(\xi)\sigma(\varphi(\xi))^{-1}\bigr)x^\nu,
 \qquad \sigma=\operatorname{Int}(s),\quad \nu=(\mu-1)/2.
\]
They do not construct a full profinite lift with leading term $q_7$
or decide this equation for such a lift if one is supplied. They also
do not establish that all pentagon solutions satisfy it.
These are the residual questions of LM2 B.m3 and E.3, anchored at
\S47, folios 416--420, scans 231--240. Their literal profinite
word evaluation has not been computed or reconstructed from the
associated graded data here.
\end{statement}
The lifts in B.12 are pro-$\ell$ lifts for $\ell\ge5$. They do not by
themselves give a full profinite word on which this $\varphi$ can be
evaluated: the quotient $v:F_0\to S_3$ does not factor through the
absolute pro-$\ell$ group for these primes. The $q_7$ certificate in
B.27 concerns its leading rational derivation.
The precise distinction is between a specified action on the source
relative completion and the existence of some peripheral graded action
on its radical. In a profinite solution the action on the open subgroup
$N$ is induced by the given action on $F_0$; it is not freely chosen
after passage to $U(F_0)$. B.26 deliberately allows that extra choice
and shows that the resulting existential condition is automatic. It
therefore cannot supply the missing implication for the prescribed
action. The absolute-completion obstruction of LM2 E.4 has been
replaced by the correctly typed rank-seven diagram, and that diagram
has been solved at the associated-graded level. The remaining work is
to specify and compute the actual relative action induced by a full
profinite word, or prove a comparison that recovers it from additional
data. Simultaneous filtered peripheral formality and passage from a
graded relative lift to such a profinite action have not been assumed.
